%% ch17 --- Taming chaos: control and synchronisation
%%
%% Written September 2026. The other edge of Part II: if a tiny disturbance
%% changes a chaotic system's future a great deal, a tiny, well-timed push can
%% steer it. Every number the reader cannot do in their head comes from
%% code/measure/ch17.py, which imports the SAME functions the listings print
%% (code/ch17/). The logistic runs and the Lorenz synchronisation use only
%% +, -, * and /, so every step count and every gap on the page is the same
%% on every machine.
%%
%% WHAT THIS CHAPTER MAY NOT SAY. It does not re-measure how fast errors grow
%% (Chapters 6, 7 and 9 own that), it does not argue whether the heart's own
%% irregularity is chaos (Chapter 15 does, and calls it contested), and it
%% does not turn the book into advice (Chapter 18's field guide does). Control
%% of chaos is stated as a laboratory result, never as a routine clinical or
%% industrial technique.
%%
%% Twin: chapters/pl/ch17-taming.tex. Section for section, macro for macro,
%% number for number; tools/parity.py compares them.

\chapter{Taming chaos: control and synchronisation}
\label{ch:taming}
\index{control of chaos}\index{synchronisation}

\noindent
Balance a broom upright on the palm of your hand. Left alone it falls:
upright is a position it cannot keep. Yet you can keep it there for as long
as your patience lasts, with movements of the hand so small that somebody
watching might not see them. You do not overpower the broom. You watch which
way it is starting to go and make a small correction, again and again.

Part~II of this book showed that a chaotic system turns a tiny difference
into a large one, and measured how fast. This chapter turns that sentence
round. If a tiny disturbance can change the future of a chaotic system a
great deal, can a tiny, well-chosen disturbance change it the way you want?
And if two chaotic systems left to themselves always drift apart, can they
be made to agree?

\begin{outcomes}
\outcome{find a fixed point of the logistic map that the rule cannot keep,
  and count how often a chaotic orbit comes near it}
\outcome{work out the small change of a parameter that cancels the miss one
  step would make}
\outcome{run a small-nudge controller on a chaotic map, and measure how long
  it waits and how large its nudges are}
\outcome{drive a copy of the Lorenz system with one number from the original
  and measure the gap between the two as time goes on}
\outcome{say where control and synchronisation of chaos have been shown to
  work, and what they are not}
\end{outcomes}

\section{A point the orbit cannot keep}
\label{sec:ch17-point}
\index{fixed point!unstable}

\noindent
Take the logistic map of Chapter~\ref{ch:logistic}, the rule that turns $x$
into $r\,x(1 - x)$, at $r = \num{3.9}$, where the orbit never settles and
never repeats. The rule still has a fixed point: the value $1 - 1/r$, which
the rule leaves exactly where it is. At $r = \num{3.9}$ that is about
\val{ch17.point}. Chapter~\ref{ch:logistic} also said why an orbit does not
stay there. Near a fixed point, one step multiplies a small miss by the slope
of the rule at that point \dash{} the stretch factor of
Chapter~\ref{ch:feedback} \dash{} and here the slope is \val{ch17.stretch}. A
miss of a thousandth becomes nearly two thousandths on the other side, then
nearly four, then nearly seven. The fixed point is a broom standing on its
end.

\pyregion{code/ch17/hidden.py}{point}{A fixed point, and what one step does to a miss there}{lst:ch17-point}

\noindent
The line that starts \code{STRETCH} asks the book's library for the slope of
the rule at the fixed point: how much one step multiplies a small miss there.

\begin{think}
The orbit at $r = \num{3.9}$ is chaotic, and the fixed point pushes away
anything that comes near it. Once the orbit has left, does it ever come back
within a hundredth of \val{ch17.point}: never, a handful of times, or again
and again?
\end{think}
\reveal{Again and again. In \num{10000} steps from $x = \num{0.2}$ it is
within a hundredth of the fixed point on \val{ch17.visits} steps, the first
time at step \val{ch17.first.visit}.}
A chaotic orbit wanders over its whole range and skips no part of it:
Chapter~\ref{ch:lorenz} found that how often a chaotic system visits each
region is one of the things you can predict about it. The listing's function
\api{visits} counts the steps that land close:

\transcript{ch17-hidden}

\noindent
Look at how the visits come: in runs of neighbouring steps. When the orbit
arrives close, the slope of \val{ch17.stretch} flings it off only gradually,
nearly doubling its miss and swapping its side at every step, so it lingers
a few steps before it is thrown clear.

The fixed point is only the simplest thing the rule cannot keep. There are
cycles of two values, of three, and of many more lengths, every one of them
unstable, and the chaotic orbit passes near each of them in turn. A chaotic
attractor is not a formless blur. It is a tangle of unstable cycles, and the
orbit hops for ever from the neighbourhood of one to the next.

\section{A nudge at the right moment}
\label{sec:ch17-nudge}
\index{Ott, Grebogi and Yorke}\index{control of chaos!nudge}

\noindent
In 1990 Edward Ott, Celso Grebogi and James Yorke published a paper called
\enquote{Controlling chaos}, with an idea as plain as the broom. Do not fight
the chaos. Wait until the orbit passes close to the unstable orbit you want,
then give the system small pushes, each only as big as it takes to cancel
the miss that the next step would make. What they pushed was a
\emph{parameter}: a setting of the system that you are allowed to adjust a
little. For the logistic map that setting is $r$.

\begin{think}
Suppose you may change $r$ at every step, but never by more than one per
cent: to anywhere between \num{3.861} and \num{3.939}. Left alone at
\num{3.9} the map is chaotic, and one per cent is hardly a change. Can
nudges that small hold the orbit at \val{ch17.point}?
\end{think}
\reveal{Yes. In the run of the next section, a controller whose nudges never
exceed one per cent is switched on at step $100$, catches the orbit at step
\val{ch17.catch}, and holds it at the fixed point for as long as the run
lasts. The largest nudge it ever makes is \val{ch17.largest.pct} per cent of
$r$.}

\begin{trapbox}
\textbf{\enquote{Chaos cannot be controlled.}} What chaos forbids is
prediction far ahead: Chapter~\ref{ch:lyapunov} measured how a small error
grows until a forecast is worthless. A controller never looks far ahead. It
looks one step ahead, and over one step a chaotic rule is as predictable as a
cup of tea: a miss is multiplied by \val{ch17.stretch}, not by a million.
Chaos takes away the long view and leaves the short one, and the short one is
all that steering needs. The chaos even does half the work, because it is
what brings the orbit near the target in the first place.
\end{trapbox}

\noindent
How big must each push be? Two facts decide it. The first you have already:
one step multiplies a miss by \val{ch17.stretch}. The second is what the knob
does. The rule is $r$ times $x(1 - x)$, so raising $r$ by a small amount $c$
for one step lifts the next value by $c$ times $x(1 - x)$. Near the fixed
point $x(1 - x)$ is about \val{ch17.lever}; call it the \emph{lever}.

\begin{think}
The orbit is \num{0.001} above the fixed point. Left alone, the next step puts
it $\num{1.9} \times \num{0.001} = \num{0.0019}$ below. Raising $r$ by $c$
lifts the next value by about \val{ch17.lever} times $c$. What $c$ cancels the
miss?
\end{think}
\reveal{$c = \num{0.0019} / \val{ch17.lever}$, about \num{0.01}: raise $r$
from \num{3.9} to \num{3.91} for this one step, about a quarter of one per
cent.}
In general the change is the miss times the slope, divided by the lever, with
its sign turned round: about \val{ch17.gain} times the miss. That is the whole
controller.

\pyregion{code/ch17/control.py}{rule}{The controller: cancel the coming miss, or wait}{lst:ch17-rule}

\noindent
\api{nudge} is the rule, and its last line is the patience. If the change it
would need is bigger than \code{largest}, it returns zero and the map runs on
untouched. With \code{largest} set to one per cent of $r$, the controller can
act only when the orbit is within \val{ch17.reach} of the fixed point, and it
has to wait for the chaos to bring it there.

\mermaidfig{ch17-wait-catch}{The controller's three moods. It waits while the
orbit is out of reach, catches it when the chaos brings it close, and holds it
with nudges; knocked away, it waits again.}{wait, catch, hold}

\section{Small change, not force}
\label{sec:ch17-small}
\index{control of chaos!cost}

\noindent
Now run the controller. The listing below leaves the map alone for $100$ steps from
$x = \num{0.2}$, then switches the controller on with nudges of at most one
per cent, and prints the state now and then.

\begin{think}
Once the orbit is caught, how big must the nudges be to keep it there: as big
as the first one, bigger because the chaos keeps pushing back, or smaller?
\end{think}
\reveal{Smaller, and then almost nothing. Within \val{ch17.settle} steps of the
catch every nudge is below a billionth of $r$, and all the nudges of the run,
added together, come to \val{ch17.total}.}

\begin{trapbox}
\textbf{\enquote{Controlling a system means overpowering it.}} The controller
supplies none of the motion. The fixed point, and the orbit's visits to it,
belong to the map; the controller only cancels the small misses it can see
coming, and a caught orbit has almost no miss left to cancel. You are
choosing which of the system's own behaviours to keep, not forcing a new one
on it. That is why Ott, Grebogi and Yorke could hold a chaotic system on any
of its many unstable cycles, and switch it between them, for the same small
nudges.
\end{trapbox}

\pyregion{code/ch17/control.py}{run}{Chaos for a hundred steps, then the controller}{lst:ch17-run}

\transcript{ch17-control}

\noindent
Python writes $-\num{2.72} \times 10^{-2}$ as \code{-2.72e-02}, which is
$-\num{0.0272}$. Read the printout downwards. For $100$ steps the orbit
wanders. The controller is switched on and for \val{ch17.waited} steps does
nothing at all, because the orbit is out of reach. At step \val{ch17.catch} it
arrives within reach, is nudged, and lands much closer; every nudge after
that is far smaller than the one before, and from step $200$ on, $x$ prints
the same ten digits.

\plotfig{ch17-control}{The logistic map at $r = \num{3.9}$, left alone for 100
steps and then controlled. The dashed line is the fixed point. Below, the
change made to $r$: zero, apart from one small dip at the catch.}{control
switched on mid-run}

The tempting alternative is force: turn the knob until the chaos stops.
Chapter~\ref{ch:logistic} found that the fixed point is stable only while $r$
is below $3$. Set $r = \num{2.9}$ and leave it there, and the orbit does
settle, at \val{ch17.brute.point}. That is a change of \val{ch17.brute.pct}
per cent of $r$, made for good, and it does not even give you the state you
wanted: the calm map's fixed point, $1 - 1/\num{2.9}$, is somewhere else. The
value \val{ch17.point} is a fixed point only when $r = \num{3.9}$, where it is
unstable, so no fixed setting of the knob can give you that state. The
controller holds it for small change: all its nudges added together come to
less than a thirtieth of the change that force makes at every single step.

Small nudges have a price, and it is paid in time. The controller can act
only when the orbit is within reach, and the reach grows in step with the
largest nudge you allow.

\begin{think}
You allow nudges ten times smaller: a tenth of one per cent of $r$. What
happens to the wait before the catch?
\end{think}
\reveal{It becomes about ten times longer. Averaged over a thousand starting
points the wait is \val{ch17.wait.big} steps with nudges of up to ten per
cent, \val{ch17.wait.mid} at one per cent and \val{ch17.wait.small} at a tenth
of one per cent.}

\pyregion{code/ch17/waiting.py}{wait}{How long until the first nudge?}{lst:ch17-wait}

\transcript{ch17-waiting}

\noindent
The reach shrinks tenfold, so the orbit has a target ten times narrower to
pass through, and on average it takes ten times as long to find it. For any
single run you cannot say in advance how long the wait will be \dash{} it is
a chaotic orbit's first visit \dash{} and the longest of the thousand runs at
one per cent waited \val{ch17.wait.longest} steps. The average, like the
statistics of Chapter~\ref{ch:lorenz}, is steady.

\begin{lab}{l17_01_controller}{A controller of your own}
Write \code{nudge} and \code{hold} for the logistic map at any $r$ at which
the fixed point $1 - 1/r$ has become unstable. At the fixed point the slope
is $r(1 - 2x)$ and the lever is $x(1 - x)$. The tests catch and hold the orbit
at $r = \num{3.9}$ and again at $r = \num{3.7}$, and check that no nudge is
ever bigger than you allowed.
\end{lab}

\section{Two chaotic systems that agree}
\label{sec:ch17-sync}
\index{synchronisation!drive and response}\index{Pecora and Carroll}

\noindent
Chapter~\ref{ch:lorenz} ran Lorenz's three equations twice, from starts a
hair apart, and watched the two runs part company for good. Two chaotic
systems left to themselves never stay together. In 1990 Louis Pecora and
Thomas Carroll asked what happens when one of them is not left to itself.
Build a copy of Lorenz's system that has its own $y$ and $z$, obeying the same
rules as the original's, but no $x$ of its own: at every moment it is handed
the original's $x$ instead.

\mermaidfig{ch17-drive}{Drive and response. The original runs freely; one of
its three numbers is passed to a copy that has only the other
two.}{one signal, from original to copy}

\begin{think}
Start the copy's $y$ and $z$ at zero, a distance of \val{ch17.sync.start} from
the original's. Both are chaotic. Will the copy's $y$ and $z$ drift away from
the original's, as the two runs of Chapter~\ref{ch:lorenz} did, or come
together?
\end{think}
\reveal{Together, and fast. The gap is below a millionth at time
\val{ch17.sync.tol}, and at time \val{ch17.sync.exact} the copy's two numbers
equal the original's in every digit the computer stores.}

\pyregion{code/ch17/sync.py}{field}{The original, and a copy handed its x}{lst:ch17-field}

\transcript{ch17-sync}

\noindent
The middle column is the copy handed $x$. The last is a second copy, from the
same wrong start, left to run on its own. The only difference between them is
that one number passed across, and it makes the whole difference: the copy
left alone never comes closer to the original than \val{ch17.free.min}, while
the driven copy's gap falls by a factor of about \val{ch17.sync.factor} with
every unit of time. Once the two agree to the last digit they stay that way,
because from then on the computer is doing the same arithmetic twice.

Why does it work? The copy is not in the loop that makes the original
chaotic. In Lorenz's system $x$ drives $y$ and $z$, and $y$ and $z$ drive $x$
back, and a small difference goes round that loop and grows. The copy only
listens. Given the same $x$, its $y$ and $z$ forget where they started, the
way two cups of tea in one room forget their difference in
Chapter~\ref{ch:models}.

\plotfig{ch17-sync}{The gap between each copy and the original, on a
logarithmic scale. Handed $x$, the copy falls into step along a straight line,
faster than the guaranteed rate; left alone, it stays as far away as the
attractor allows.}{a copy falling into step}

\begin{trapbox}
\textbf{\enquote{Two chaotic systems can never agree.}} Two chaotic systems
left to themselves can never stay together, because each amplifies whatever
difference there is between them. Coupled so that one listens to the other,
they can agree exactly, chaos and all. Notice what the agreement is not: the
copy agrees with the original now, not with the original's future.
Synchronisation is not prediction.
\end{trapbox}

\begin{rigourbox}
Why it must work, whatever $x$ does. Write $e_y$ and $e_z$ for how far the
copy's $y$ and $z$ are from the original's. Subtract the copy's rules from the
original's and work out how fast $e_y^2 + e_z^2$ changes: every term that
contains $x$ cancels, and what is left, $-2e_y^2 - 2\beta e_z^2$ with
$\beta = 8/3$, is never positive. From that one line it follows that the gap
shrinks at least by the factor $e$, about \num{2.718}, in every unit of time;
the measured factor of \val{ch17.sync.factor} is better than the guarantee.
An argument of this kind is in Strogatz's textbook \emph{Nonlinear Dynamics
and Chaos}. Not every choice of the number handed across works, and Pecora
and Carroll's paper gives the test that tells.
\end{rigourbox}

\begin{lab}{l17_02_sync}{How long does the copy take?}
Write \code{copy\_field}, the five rates of change of the original and the
copy; \code{error}, the distance between the copy's $y$ and $z$ and the
original's; and \code{sync\_time}, which steps the five numbers together and
returns the first time the error falls below a tolerance. The tests check
that three more factors of ten take no longer than the guarantee above
allows.
\end{lab}

\section{What has been done with it, and what has not}
\label{sec:ch17-world}
\index{secure communication}\index{control of chaos!in the laboratory}

\noindent
Pecora and Carroll's result led at once to a proposal for keeping secrets.
Add a faint message to the original's $x$ and send the sum. To anyone
listening it sounds like noise. A receiver built as the matching copy falls
into step with the chaotic part of the signal, and subtracting what it
rebuilds leaves the message. Circuits that do this were built soon afterwards.
Later work showed that many such schemes can be broken by an eavesdropper who
studies the signal, and chaos has not replaced ordinary encryption.

\begin{worldbox}
Within a few years of 1990 the small-nudge method had held chaotic lasers and
electronic circuits on orbits they would not keep by themselves. In 1992 Alan
Garfinkel and colleagues applied it in an experiment on isolated rabbit heart
tissue, whose beating had been made irregular with a drug: electrical stimuli,
timed by the same wait-then-nudge logic, made the rhythm regular again. Space
missions use the same double edge of sensitivity. Near the places where the
pulls of the Sun, the Earth and the Moon nearly balance, the three-body
problem of Chapter~\ref{ch:mechanics} is chaotic, and a small burn of the
engines changes a path a great deal. NASA's ISEE-3 spacecraft was steered
away from its station between the Sun and the Earth in this way, with a
series of close passes of the Moon and little fuel, and in 1985 it flew
through the tail of a comet.
\end{worldbox}

\begin{warning}
Keep these in proportion. Control of chaos has been demonstrated many times in
the laboratory. It is not a routine clinical treatment, and it is not a
standard industrial technique; whether the heart's own irregularity is chaos
at all is a question that Chapter~\ref{ch:life} calls contested.
\end{warning}

\noindent
What this chapter adds to the book's argument is this. Part~II measured
sensitivity as a limit: it is why the horizon of a forecast is short. It is
also a lever. A calm system has one behaviour, and changing it takes force. A
chaotic system is full of behaviours it cannot keep, and any of them can be
chosen, and held, for small change, as long as you are willing to watch and
wait. Chaos is not only a limit on what you can know. It is a resource.
Chapter~\ref{ch:why} gathers what the whole book has measured into a guide
for living with both.

\begin{summarybox}
\begin{itemize}
\item At $r = \num{3.9}$ the logistic map has a \textbf{fixed point it cannot
  keep}: one step multiplies a miss by the slope, \val{ch17.stretch}. A
  chaotic orbit keeps passing close to it, and to many \textbf{unstable
  cycles}.
\item A \textbf{controller} waits for a close pass, then changes the
  \textbf{parameter} $r$ by the amount that cancels the coming miss: the miss
  times the slope, divided by the \textbf{lever}, with the sign turned round.
  It predicts only one step ahead, and chaos allows that.
\item Once the orbit is caught the nudges shrink to almost nothing. Force, a
  large and permanent change of $r$, cannot even give the same state. Smaller
  nudges cost a longer \textbf{wait}, whose average is steady.
\item A copy of the Lorenz system handed the original's $x$ falls into step
  with it: \textbf{synchronisation}. Left alone it never does. Agreement is
  not prediction.
\item Control and synchronisation of chaos are \textbf{laboratory results},
  in lasers, circuits and heart tissue, and a tool of spacecraft steering.
  Sensitivity is a \textbf{resource} as well as a limit.
\end{itemize}
\end{summarybox}

\canyou

\begin{checkpoint}
\item For the logistic map at $r = \num{3.5}$, what is the fixed point
  $1 - 1/r$, what is the slope $2 - r$ there, and can an orbit stay at it?
  \answerto{About \num{0.714}, with slope $-\num{1.5}$. A miss is multiplied
  by \num{1.5} and flips side at every step, so the orbit cannot stay.}
\item At $r = \num{3.9}$ the orbit is \num{0.002} below the fixed point.
  Where does one uncontrolled step put it, and what change of $r$ cancels the
  miss?
  \answerto{About \num{0.0038} above it. Lower $r$ by about
  $\num{0.0038} / \val{ch17.lever}$, that is by \num{0.02}, for that one step.}
\item Why can a controller steer a chaotic system when nobody can forecast
  it far ahead?
  \answerto{Because the controller needs to predict only one step ahead, and
  over one step a miss grows by a factor of about two, not exponentially
  without end.}
\item You halve the largest nudge you allow. What happens to the reach of the
  controller, and roughly to the average wait?
  \answerto{The reach halves, so the target is half as wide and the average
  wait roughly doubles.}
\item Why does setting $r = \num{2.9}$ for good not give you the state
  \val{ch17.point}?
  \answerto{Because the calm map settles at its own fixed point,
  $1 - 1/\num{2.9}$, about \val{ch17.brute.point}. The value \val{ch17.point}
  is a fixed point only at $r = \num{3.9}$, where it is unstable.}
\item In Pecora and Carroll's arrangement, what is passed from the original
  to the copy, and does the copy then predict the original's future?
  \answerto{The original's $x$. No: the copy agrees with the original at each
  moment, which is synchronisation, not prediction.}
\item Name two places where control of chaos has been demonstrated, and say
  what it is not.
  \answerto{Any two of lasers, electronic circuits and isolated heart tissue
  in the laboratory. It is not a routine clinical or industrial technique.}
\end{checkpoint}
